% ============================================================
% OP_GUT1_zeromode_programme_v1.tex  (PROPOSAL-ONLY)
% Sharpening of the colour zero-mode problem: axioms, admissible
% backgrounds, route taxonomy, precise sub-problems, failure
% modes. NO candidate operator constructed, NO existence claim,
% NO preprint edits by this file.
% ============================================================
\section*{OP-GUT1 zero-mode programme (proposal-only)}

\subsection*{0. Scope}
This document sharpens the block ``Definition (the colour
zero-mode problem) / Open problem (OP-GUT1, sharpened)'' of
\S colour\_completion into a workable programme: precise axioms
for $\mathcal D_{\rm defect}$, the admissible background classes
already present in the ECT defect ontology, a taxonomy of three
candidate routes with honest assessments, and exact formulations
of the induced-$(h,\Omega)$ sub-problems.
Everything below is Open; nothing upgrades P7.

\subsection*{1. Axioms for $\mathcal D_{\rm defect}$ (sharpened)}
\begin{enumerate}
  \item Locality: a differential operator built pointwise from the
    minimal condensate data $(\rho,\theta,n^A)$ and their
    derivatives; no external structure.
  \item Ellipticity, first-order Dirac type preferred (kernel
    robustness, index tools).
  \item Equivariance under the residual ordered-branch symmetry
    $O(3)\times U(1)_\theta$.
  \item A fixed $U(1)_\theta$ charge assignment on the domain
    bundle: the $\theta$-action is the natural candidate for the
    complex structure on the kernel (multiplication by~$i$), since
    a \emph{complex} rank-3 module is required, not a real one.
  \item Spectral gap separating the would-be zero modes from the
    continuum on the relevant backgrounds, so that
    $\ker\mathcal D_{\rm defect}$ forms a smooth bundle over
    spacetime and the Berry--Wilczek--Zee connection is defined
    (no level crossing along physical paths).
\end{enumerate}

\subsection*{2. Admissible backgrounds (grounded in the existing
ontology)}
Per \S ``Types of defect and topological configurations'' and
app.~relic\_homotopy:
(i)~$\theta$-winding / vortex-core configurations (compact phase
fibre; modulus suppression at the core);
(ii)~$\pi_3(S^3)=\mathbb Z$ texture/skyrmion sectors --- the only
nontrivial homotopy class of the primary $O(4)\to O(3)$ breaking
($\pi_1=\pi_2=0$: no strings, no monopoles);
(iii)~domain-wall-like / interface configurations, including
ordered--Euclidean interfaces at the BR1 boundary;
(iv)~composite configurations.
The ontology itself is marked ``not yet fully classified''; the
programme inherits that conditionality explicitly.

\subsection*{3. Route taxonomy (no commitments)}
\textbf{R1 --- index/charge route ($\theta$-winding
backgrounds).}
Dirac-type operators coupled to winding-$n$ vortex backgrounds
generically have $\dim\ker=|n|$ by Jackiw--Rossi/Callias-type
counting (references to be page-verified before any use:
Jackiw--Rossi NPB 190 (1981) 681; Callias CMP 62 (1978) 213;
E.~Weinberg PRD 24 (1981) 2669).
Honest assessment: this route \emph{transforms} the rank problem
rather than solving it --- $\dim_{\mathbb C}\ker=3$ becomes ``why
are $|n|=3$ cores the relevant stable objects'', a
charge-selection question (energetics of binding three elementary
cores?) that is itself entirely open.
Recordable as a structural observation; no claim.

\textbf{R2 --- representation/degeneracy route (residual-symmetry
multiplet).}
The kernel carries an irreducible multiplet of the residual
$O(3)$; the three-dimensional vector representation would give
$\dim_{\mathbb R}\ker=3$, complexified to $\mathbb C^3$ by the
$U(1)_\theta$ action (axiom~4).
Two striking structural matches, recorded as \emph{questions},
not results:
(a)~the antisymmetric invariant $\varepsilon_{abc}$ of P7 would
be the $\varepsilon_{ijk}$ of the residual $O(3)$ vector triple
product, i.e.\ $\Omega$ induced by rotational geometry;
(b)~$N_c=3$ would tie to $d_{\rm spatial}=3$.
Honest flag: this is the most ECT-native candidate and also the
most seductive (numerological risk); the existence of such a
triplet kernel is completely open and must never be asserted.

\textbf{R3 --- interface route.}
Localised modes on ordered--Euclidean or inter-branch interfaces
(Jackiw--Rebbi-type localisation); counting data unclear; least
developed; listed for completeness.

\textbf{R4 --- $G_2$ reduction.}
Nested deeper candidate ($u^\perp\simeq\mathbb C^3$); deferred to
the separate $G_2$ mini-round by review decision; listed only for
completeness of the taxonomy.

\subsection*{4. The induced-$(h,\Omega)$ sub-problems, made
precise}
\textbf{(h).} For elliptic $\mathcal D_{\rm defect}$ the
fibrewise $L^2$ inner product on $\ker$ supplies $h$
automatically; the substantive content is smoothness of the
kernel bundle (axiom~5), not the metric.
\textbf{($\Omega$).} The $U(3)\to SU(3)$ reduction is equivalent
to a canonical trivialisation of the complex determinant line of
the kernel bundle:
\[
\Omega\ \text{exists canonically}
\iff
c_1\bigl(\det\nolimits_{\mathbb C}\ker\mathcal D_{\rm defect}\bigr)=0
\ \text{with a preferred section from condensate orientation data.}
\]
This turns the ``why $SU(3)$ rather than $U(3)$'' question into a
precise index-bundle condition.
Route-specific candidates: in R2, $\Omega=\varepsilon_{ijk}$ from
the residual $O(3)$; in R1, only a non-canonical ordering of
winding modes --- a recorded weakness of R1.
\textbf{(BWZ).} The colour connection is the Berry--Wilczek--Zee
connection of the smooth kernel bundle; its Yang--Mills kinetics
is \emph{not} part of this programme (separate question, partially
addressed by the minimal colour EFT section).

\subsection*{5. Failure modes (recorded honestly)}
Generic $\dim\ker=0$ on natural backgrounds; kernel real with no
$\theta$-complexification; $|n|=3$ selection unjustified (R1);
level crossing destroying the bundle; $c_1\neq0$ obstruction to
$\Omega$; and even full success would leave Yang--Mills kinetics,
$g_s$, confinement, and OP-GUT2 representations open.
Consistency check inherited from the $N_c$ classification: any
successful mechanism must produce a \emph{complex} triplet
(chiral-anomaly-relevant), and the classification guarantees that
no anomaly argument can substitute for it.

\subsection*{6. Deliverables if approved (future round; nothing
executed now)}
D1 (small): a ``Candidate routes (all Open)'' paragraph after the
``Open problem (OP-GUT1, sharpened)'' block: R1--R3 in three
sentences each, plus the two precise sub-statements (R1
rank$\to$charge transformation; $\Omega\iff$ determinant-line
trivialisation).
D2 (optional): bib additions JackiwRossi1981 / Callias1978 /
WeinbergE1981 after page-level verification.
D3: statuses untouched everywhere; P7 remains a postulate;
OP-GUT1 remains Open/Reframed.

\subsection*{7. Questions for GPT}
(Q1)~Approve the route taxonomy R1--R3 (R4 deferred) for eventual
publication shape?
(Q2)~Is the R2 double observation
($\Omega\leftrightarrow\varepsilon_{ijk}$ of residual $O(3)$;
$N_c=3\leftrightarrow d_{\rm spatial}=3$) admissible in the text
as an explicitly flagged open question with firewalls, or too
seductive to print at all?
(Q3)~Include the determinant-line/$c_1$ reformulation of
$U(3)\to SU(3)$ as a precise sub-problem statement?
(Q4)~Include the R1 reformulation (rank selection $\mapsto$
charge-3 selection, energetics open) --- useful honesty or
unnecessary speculation?
(Q5)~Reference set for D2: Jackiw--Rossi + Callias + E.~Weinberg,
or a subset?
